\documentclass[a4paper]{article}
\usepackage[pdftex]{graphicx}
\usepackage[utf8]{inputenc}
\usepackage{enumerate}
\usepackage{icomma}
\usepackage{amssymb}
\usepackage{tikz}
\usepackage{href-ul}
\hypersetup{
	colorlinks=true,
	linkcolor=blue,
	urlcolor=blue}
\usepackage{geometry}
\geometry{a4paper, top=15mm, left=15mm, right=15mm, bottom=15mm,
	headsep=10mm, footskip=12mm}

\begin{document}
\begin{tikzpicture}(10,3)
	\draw[ultra thick](2,0) --(10,0) -- (10,1.3) --(2,1.3) -- (2,0);
	\draw[fill=black](2,0)-- (0,.6) -- (2,.6) -- (2,0);
	\node at (6,.6) {\large Lösungsblatt von \href{https://www.okuyakl.de}{www.okuyakl.de}};
\end{tikzpicture}
\vspace{0.5 cm}

\noindent{\bf Aufgabe 1.}\\
\renewcommand{\arraystretch}{2}
\begin{tabular}{rlp{20 pt}rl}
	a) & $\infty$ & & b) & $\infty$ \\
	c) & 4 & & d) & 0 \\
	e) & 0 & & f) & ex. nicht \\
	g) & $\infty$ & & h) & 0 
\end{tabular}

\noindent{\bf Aufgabe 2.}\\

$$\lim\limits_{x \to \infty} {3x-5 \over 4x+3}= 
\lim\limits_{x \to \infty} \frac{x(3 - {5 \over x})}{x(4 - {3\over x})}= 
{3 -0 \over 4 -0}={3 \over 4} $$
$$\lim\limits_{x \to -\infty} {4x+x^3 \over x^4-2}= 
\lim\limits_{x \to -\infty} \frac{x^4({4\over x^3} + {1 \over x})}{x^4(1 - {2\over x^4})}=
{0 +0\over 1+0} =0$$

\noindent{\bf Aufgabe 3. a)}\\
Direkt 0 einsetzen:
$$= -{5 \over 2} $$
\noindent{\bf Aufgabe 3. b)}\\
$x$ im Nenner ausklammern und kürzen, dann 0 einsetzen:
$$=\lim\limits_{x \to 0}{5x^2 \over x-1} = {0 \over -1} = 0 $$
\noindent{\bf Aufgabe 3. c)}\\
Im Zähler 2 einsetzen, im Nenner Vorzeichen für Näherung beachten:
$$ "{1+4 \over 2 \cdot (+0)}"= + \infty$$
\noindent{\bf Aufgabe 3. d)}\\
Zählergrad=Nennergrad; Faktoren beachten:
$$= {-3 \over 2}$$
\noindent{\bf Aufgabe 3. e)}\\
Faktorisieren, kürzen und einsetzen:
$$=\lim\limits_{x \to 4}{3(x-4)\over (x+5)(x-4)}=\lim\limits_{x \to 4} {3 \over x+5} = {1 \over 3}$$
\noindent{\bf Aufgabe 3. f)}\\
Zählergrad $>$ Nennergrad; Ausklammern, kürzen, einsetzen:
$$= \lim\limits_{x \to -\infty}{x(7x-{2\over x})\over x(2+{3\over x})}= \lim\limits_{x \to -\infty}{7x-{2\over x}\over 2+{3\over x}}={-\infty \over 2}=-\infty $$

\begin{center}
	\includegraphics[width=7 cm]{../../viecher/endcomic.pdf}
	
	Hier geht es zurück zum \href{https://www.okuyakl.de/math/m10limL034/aa034.pdf}{Aufgabenblatt}
\end{center}

\end{document}

